\subsection{Sender Side (RTC Transmission Pipeline)}

\begin{figure}[H]
    \vspace{1em} 
    \centering
    \includegraphics[width=1\textwidth]{images/methodology/sender_side.png}
    \vspace{0.25cm}
    \caption{Sender-side RTC Transmission Pipeline}
    \label{fig:sender_side}
\end{figure}

The sender-side of RTC transmission pipeline is responsible for transforming a speaker's voice into a compressed network stream suitable for real-time communication. In the proposed framework, the sender-side pipeline is treated as a standard RTC subsystem and is not modified by the detection mechanism and serves as a realistic source of audio data that reflects practical deployment conditions found in VoIP and conferencing systems.

Modeling the sender-side pipeline is important because audio transmitted through RTC systems undergoes multiple transformations before reaching the receiver. These transformations introduce compression artifacts, quantization effects, and network-induced distortions that may affect the discriminative characteristics used for deepfake detection. Consequently, the proposed system evaluates detection performance under realistic communication constraints rather than on pristine studio-quality recordings.


\subsubsection{Audio Capture and Pre-processing}

The transmission process begins with acquisition of speech through a microphone device. The analog speech waveform is converted into a digital signal through Analog-to-Digital Conversion (ADC), where sampling and quantization are performed at a predefined sampling rate (typically ranging from 16 kHz to 48 kHz).

Let the continuous-time speech signal be represented as $x(t)$ which is transformed into a discrete-time representation $x[n]$ through the sampling process.

Depending on the operating system and device capabilities, additional preprocessing may be applied, including gain normalization, noise suppression, and acoustic echo cancellation. These operations improve speech intelligibility but may also alter spectral characteristics relevant to deepfake analysis.

Here, the signal remains in a near-natural form but may already include environmental noise and recording artifacts. This stage introduces uncontrolled variability (speaker device, environment), which contributes to domain diversity in downstream detection.

\subsubsection{Encoding and Compression}

The preprocessed audio signal is compressed using a speech codec before transmission. Modern RTC systems typically employ codecs such as Opus due to their low latency and adaptive bitrate capabilities. 

Let $x[n]$ denote the input waveform and $C(\cdot)$ denote the codec transformation:
\begin{equation}
x_c = C(x[n])
\label{eq:codec_transform}
\end{equation}
where $x_c$ is the compressed bitstream representation.

Compression reduces transmission bandwidth requirements but introduces codec-specific distortions, including quantization noise, spectral smoothing, and loss of fine-grained acoustic information. These effects may partially mask or modify artifacts generated by synthetic speech systems.

\subsubsection{Packetization and Stream Generation}

The compressed audio stream is segmented into small frames suitable for real-time delivery. Each frame is encapsulated into transport packets containing timing and synchronization metadata.

The packetization process is represented as
\begin{equation}
x_c \rightarrow \{p_1, p_2, \ldots, p_n\}
\label{eq:packetization}
\end{equation}
where $p_i$ represents an individual network packet.

Each packet generally contains a compressed payload along with sequence numbers and timestamps required for stream reconstruction at the receiver.